\paragraph{New payment terms leave earlier assessments intact.}
A prospective amendment answers a different question from a corrected source
reading: what changes from the amendment date onward? Return to the earlier
single-invoice example, before the illustrative second sale. The original
invoice was due on 31~January 2027. Its remaining USD~20 was assessed as
overdue on 28~February. Suppose the company and customer are now modeled as
extending its due date on 10~March to 30~April. The February assessment remains
USD~20 overdue under the original terms. A fresh assessment on 31~March records
zero overdue under the extension. Replacing the old due date in place would
lose the information needed to explain why both assessments are correct in
their respective histories.

The bounded \texttt{ReceivableTermHistory/1} reference supplies this separate
prospective operation. Its migration starts from a valid legacy claim book
and an explicit premise that no earlier amendment occurred. On that premise,
it imports the existing due date as the initial term and preserves the last
stored assessment. It does not reconstruct earlier assessments. The imported
assessment's event identity is recorded as unknown because the old summary
does not identify that event. The source book and its hash remain available;
the migration leaves its balances, cash flows, claim identities and original
scalar observations unchanged. Consequently, even a successful migration
leaves completeness of the earlier term history unresolved.

Each claim then has an append-only map of term versions and an append-only
map of assessments. A term names its due date, effective date, originating
event, predecessor and supplied premise. The current-term reference identifies
the version used for a new assessment. Each assessment retains its own term
reference, assessment date, remaining amount at that date and overdue amount.
The validator checks the term applicable at that historical date; it does not
recompute every old observation using the current due date or current balance.
The familiar scalar assessment fields remain the summary of the latest
assessment, while the current remaining balance can change independently.

For example, a second extension on 10~April moves the due date to 31~May.
A receipt of USD~5 on 15~April reduces the claim to USD~15, while the March
assessment still records the USD~20 that remained when it was made. The
30~April assessment records zero overdue; the 30~June assessment records
USD~15 overdue. A 31~December assessment retains that amount. Year-end carry
preserves the three term versions and all assessment records, and opens a new
assessment window. Collection of the final USD~15 on 10~January 2028 clears
the claim; the 31~January assessment records zero. Final cash is USD~130,
receivables are zero and retained earnings are USD~20, as in the unamended
single-invoice example: these date changes altered the assessment history,
while the supplied sale and collection amounts determined the monetary result.

The immediate amendment takes effect on its event date. It requires the
current state hash, existing claim and customer identities, the expected
current term, a new term identity and explicit typed inputs. It refuses
backdating, a due date before effectiveness, duplicate identities and two
terms with the same effective date. An assessment seals its day against later
amendments or collections. An amendment followed by assessment on that same
day is allowed, so the order is explicit. Unknown required input remains
unresolved. The bounded maps admit at most 32~terms and 64~assessments per
claim; this reference does not silently discard old records when those limits
are reached.

These constructed date-only amendments and assessments produce no journal
entry. That is the declared scope of this method; an actual contract change
may require a separately specified measurement and accounting treatment.
Legal enforceability, modification gains or losses, interest and impairment
remain outside this reference. The exact typed requests, historical states,
receipts, carry and negative cases are retained for two account charts in
\path{data/accounting-master/receivable-term-history-20261025/}. They connect
the prospective operation to the preceding source-correction distinction
without treating either operation as evidence for a real invoice population.

